\section{问题三的建模与求解}
\label{sec:problem3}
\subsection{模型建立}
\subsubsection{扰动设置}
以碳排因子或材料单价为待分析参数，在基准值附近设置相对扰动，其他条件保持一致。
参数范围及选择依据应根据实际数据补充。

\subsection{求解与检验}
\subsubsection{结果检验}
对各扰动情景重新求解，比较目标值和折中方案的变化。
在此填写敏感性分析结果，并说明稳定性检验的依据与适用范围。
